\documentclass[11pt,a4paper]{article}
\usepackage[margin=2.5cm]{geometry}
\usepackage{hyperref}
\usepackage{booktabs}
\usepackage{enumitem}
\usepackage{xcolor}
\usepackage{titlesec}
\usepackage{array}
\usepackage{parskip}

\definecolor{accent}{HTML}{00D4FF}
\definecolor{dim}{HTML}{555577}
\definecolor{warn}{HTML}{FFAA44}
\definecolor{ok}{HTML}{44FFAA}

\hypersetup{colorlinks=true, linkcolor=accent, urlcolor=accent}

\titleformat{\section}{\large\bfseries\color{accent}}{}{0em}{}[\titlerule]
\titleformat{\subsection}{\normalsize\bfseries}{}{0em}{}

\title{\textbf{VSEPR-SIM --- Graphing Module Plan}\\
       \large Python Visualisation Suite for Atomistic Simulation Data\\[0.4em]
       \normalsize\color{dim}Planning Document v1.0 --- April 2026}
\author{VSEPR-SIM V4 Research Platform}
\date{}

\begin{document}
\maketitle
\tableofcontents
\vspace{1em}
{\small\itshape\color{dim}%
This document describes the planned Python graphing module suite.
No implementation code is written here.
All modules consume JSON snapshots produced by \texttt{crystal\_card\_popup.py}
and future simulation output files.
Code will be written in a separate phase after this design is reviewed.}

% ─────────────────────────────────────────────────────────────────────────────
\section{\S{}G0 --- Design Philosophy}

The graphing suite follows the same anti-black-box principle as the rest of
VSEPR-SIM\@: every curve, every data point, and every axis must be directly
traceable to a specific simulation output.
No synthetic or interpolated data should be presented without explicit labelling.

Three core ideas govern the design:

\begin{enumerate}[leftmargin=2em]
  \item \textbf{Data loads visibly.}
        Curves animate from left to right as data arrives.
        The user sees the process, not just the result.
  \item \textbf{2D and 3D are co-equal.}
        Every dataset has a designated 2D representation (publication-ready)
        and a designated 3D representation (exploratory, rotatable).
  \item \textbf{Modules are single-purpose.}
        Each module reads exactly one data type, produces exactly one figure class.
        No omnibus dashboards at this stage.
\end{enumerate}

% ─────────────────────────────────────────────────────────────────────────────
\section{\S{}G1 --- Data Sources}

\subsection{G1.1 Card Snapshot JSON}
Produced by \texttt{crystal\_card\_popup.py} on every card open.
Location: \texttt{data/card\_snapshots/\{symbol\}\_\{lattice\}\_\{ts\_ms\}.json}

Fields available for graphing:

\begin{center}
\begin{tabular}{lll}
\toprule
Field & Type & Use \\
\midrule
\texttt{timestamp\_ms}       & int   & timeline axis \\
\texttt{symbol}              & str   & per-element colour coding \\
\texttt{lattice}             & str   & group/facet key \\
\texttt{a0\_angstrom}        & float & EOS / lattice parameter plot \\
\texttt{Tmelt\_K}            & float & thermal curve overlay \\
\texttt{Ecoh\_eV}            & float & cohesive energy chart \\
\texttt{gamma}               & float & score distribution / radar \\
\texttt{q\_data}             & float & score distribution / radar \\
\texttt{compact}             & float & score distribution / radar \\
\texttt{render\_mode}        & str   & 2D vs.\ 3D session split \\
\texttt{complexity\_label}   & str   & supercell tier histogram \\
\texttt{supercell\_dims}     & list  & atom-count scaling curve \\
\texttt{view\_azimuth\_deg}  & float & viewing angle distribution \\
\texttt{view\_elevation\_deg}& float & viewing angle distribution \\
\bottomrule
\end{tabular}
\end{center}

\subsection{G1.2 Future Simulation Output Files}
Planned formats (not yet implemented):

\begin{itemize}[leftmargin=2em]
  \item \texttt{data/eos/\{symbol\}\_eos.csv} ---
        columns: \texttt{volume\_AA3, energy\_eV, pressure\_GPa}
  \item \texttt{data/rdf/\{symbol\}\_rdf.csv} ---
        columns: \texttt{r\_AA, g\_r}
  \item \texttt{data/trajectory/\{symbol\}\_md.xyz} ---
        XYZA format; existing reader in \texttt{src/io/xyz\_format.cpp}
  \item \texttt{data/convergence/\{symbol\}\_conv.csv} ---
        columns: \texttt{step, energy\_eV, force\_max\_eVAA}
  \item \texttt{data/thermo/\{symbol\}\_thermo.csv} ---
        columns: \texttt{step, T\_K, P\_GPa, E\_total\_eV}
  \item \texttt{data/phonon/\{symbol\}\_dos.csv} ---
        columns: \texttt{freq\_THz, dos}
\end{itemize}

% ─────────────────────────────────────────────────────────────────────────────
\section{\S{}G2 --- Planned Module List}

\begin{center}
\begin{tabular}{lllp{5.2cm}}
\toprule
Module file & Type & Data source & Description \\
\midrule
\texttt{plot\_score\_dist.py}    & 2D & snapshots     & Histogram + KDE of $\gamma$, $Q$, $C$
                                                         across all recorded sessions \\
\texttt{plot\_score\_radar.py}   & 2D & snapshots     & Animated radar polygon building up
                                                         per-element as snapshots are scanned \\
\texttt{plot\_eos.py}            & 2D & eos CSV       & Energy--volume curves with
                                                         Birch--Murnaghan fit overlay \\
\texttt{plot\_rdf.py}            & 2D & rdf CSV       & $g(r)$ with animated trace
                                                         revealing from $r=0$ \\
\texttt{plot\_rdf\_3d.py}        & 3D & rdf CSV       & $g(r,\theta)$ surface in
                                                         Matplotlib 3D axes \\
\texttt{plot\_convergence.py}    & 2D & conv CSV      & Energy and force vs.\ step;
                                                         dual-axis; animated left-to-right \\
\texttt{plot\_thermo.py}         & 2D & thermo CSV    & $T$, $P$, $E$ vs.\ MD step;
                                                         triple panel; animated reveal \\
\texttt{plot\_thermo\_3d.py}     & 3D & thermo CSV    & $T$--$P$--$E$ phase-space trajectory;
                                                         points appear sequentially \\
\texttt{plot\_trajectory\_2d.py} & 2D & XYZA          & Per-atom displacement vs.\ step
                                                         as translucent line bundle \\
\texttt{plot\_trajectory\_3d.py} & 3D & XYZA          & 3D scatter of atom positions;
                                                         animated by frame \\
\texttt{plot\_lattice\_param.py} & 2D & snapshots     & $a_0$ vs.\ $T_\text{melt}$ scatter,
                                                         colour-coded by lattice type \\
\texttt{plot\_ecoh\_bar.py}      & 2D & snapshots     & $E_\text{coh}$ bar chart; bars
                                                         grow left-to-right \\
\texttt{plot\_phonon\_dos.py}    & 2D & phonon CSV    & Phonon DOS; shaded area fills in
                                                         from low to high frequency \\
\texttt{plot\_session\_timeline.py} & 2D & snapshots  & Dot-strip timeline coloured
                                                         by element; renders live \\
\texttt{plot\_supercell\_scale.py}  & 2D & snapshots  & Atom count vs.\ supercell tier;
                                                         FCC/BCC/HCP side by side \\
\bottomrule
\end{tabular}
\end{center}

% ─────────────────────────────────────────────────────────────────────────────
\section{\S{}G3 --- Animated Data-Loading Pattern}

Every module that touches time-series or step-indexed data \emph{must} implement
the animated loading pattern.
The intended visual effect: the curve or scatter plot is empty at $t=0$;
data points appear one-by-one (or in small batches) from left to right,
building the full picture in under 3 seconds.

\subsection{G3.1 Planned mechanism}
\begin{itemize}[leftmargin=2em]
  \item Load the full data array into memory first (no streaming required).
  \item Use Matplotlib \texttt{FuncAnimation} with a frame count equal to
        the number of data points (or a fixed 60--120 frames with sub-sampling).
  \item Each frame reveals the next batch of $k$ points
        ($k = \lceil N / 90 \rceil$ for approximately 1.5\,s at 60\,fps).
  \item For line plots: update \texttt{line.set\_data(x[:frame], y[:frame])}.
  \item For scatter plots: replace the \texttt{PathCollection} offsets each frame.
  \item For area/fill plots: redraw \texttt{fill\_between} on a clipped slice.
\end{itemize}

\subsection{G3.2 Axes style}
All modules share a common dark style consistent with the card UI:

\begin{center}
\begin{tabular}{ll}
\toprule
Element & Value \\
\midrule
Figure background  & \texttt{\#0d0d1a} \\
Axes background    & \texttt{\#111130} \\
Tick / label colour & \texttt{\#888888} \\
Grid colour        & \texttt{\#222244}, dashed \\
Accent colour      & \texttt{\#00d4ff} \\
Secondary colour   & \texttt{\#44ffaa} \\
Warning colour     & \texttt{\#ffaa44} \\
Font               & Consolas or DejaVu Sans Mono \\
\bottomrule
\end{tabular}
\end{center}

\subsection{G3.3 Export}
Every module saves the completed figure to
\texttt{data/figures/\{module\_stem\}\_\{timestamp\}.png}
at 150\,dpi by default, 300\,dpi with \texttt{--hires} flag.
Animated versions save \texttt{.gif} alongside.

% ─────────────────────────────────────────────────────────────────────────────
\section{\S{}G4 --- 3D Module Specifics}

3D modules use \texttt{matplotlib.pyplot} with \texttt{projection='3d'}.
No external OpenGL or VTK dependency is introduced at this stage.

\subsection{G4.1 Rotation}
All 3D plots are initialised with:
\texttt{ax.view\_init(elev=22, azim=30)} ---
the same hardcoded angles used by the crystal card viewer.

\subsection{G4.2 Depth cues}
\begin{itemize}[leftmargin=2em]
  \item Scatter point size scales with depth (nearer = larger).
  \item Alpha channel encodes depth (nearer = more opaque).
  \item Line bundles use \texttt{LineCollection} with per-segment alpha.
\end{itemize}

\subsection{G4.3 Planned 3D figures and axes}

\begin{center}
\begin{tabular}{llll}
\toprule
Module & X axis & Y axis & Z axis \\
\midrule
\texttt{plot\_rdf\_3d.py}        & $r$ (\AA)  & $\theta$ (rad) & $g(r,\theta)$ \\
\texttt{plot\_thermo\_3d.py}     & $T$ (K)    & $P$ (GPa)      & $E$ (eV)      \\
\texttt{plot\_trajectory\_3d.py} & $x$ (\AA)  & $y$ (\AA)      & $z$ (\AA)     \\
\bottomrule
\end{tabular}
\end{center}

% ─────────────────────────────────────────────────────────────────────────────
\section{\S{}G5 --- Module Directory Layout}

\begin{verbatim}
tools/
  graphing/
    __init__.py                  (empty, marks as package)
    _style.py                    (shared dark style + colour palette)
    _loader.py                   (JSON snapshot reader + CSV reader helpers)
    _anim.py                     (FuncAnimation wrapper + export helper)
    plot_score_dist.py
    plot_score_radar.py
    plot_eos.py
    plot_rdf.py
    plot_rdf_3d.py
    plot_convergence.py
    plot_thermo.py
    plot_thermo_3d.py
    plot_trajectory_2d.py
    plot_trajectory_3d.py
    plot_lattice_param.py
    plot_ecoh_bar.py
    plot_phonon_dos.py
    plot_session_timeline.py
    plot_supercell_scale.py

data/
  card_snapshots/    (JSON -- written by crystal_card_popup.py)
  eos/               (CSV  -- future EOS engine)
  rdf/               (CSV  -- future RDF engine)
  convergence/       (CSV  -- future SCF/MD engine)
  thermo/            (CSV  -- future MD thermostat)
  trajectory/        (XYZA -- future MD engine)
  phonon/            (CSV  -- future phonon engine)
  figures/           (PNG + GIF output from graphing modules)
\end{verbatim}

% ─────────────────────────────────────────────────────────────────────────────
\section{\S{}G6 --- CLI Interface (Planned)}

Each module is callable as a standalone script:

\begin{verbatim}
python tools/graphing/plot_score_dist.py [--snapshot-dir PATH] [--hires]
python tools/graphing/plot_eos.py Au Fe Ti [--overlay] [--hires]
python tools/graphing/plot_convergence.py Au [--no-anim]
python tools/graphing/plot_trajectory_3d.py Au [--frame N]
\end{verbatim}

All modules accept \texttt{--no-anim} to skip animation and render the final
state immediately (for batch or CI use).

% ─────────────────────────────────────────────────────────────────────────────
\section{\S{}G7 --- Dependencies (Minimal)}

\begin{center}
\begin{tabular}{lll}
\toprule
Library & Version & Use \\
\midrule
\texttt{matplotlib} & $\geq$\,3.7  & all 2D and 3D plots, animation \\
\texttt{numpy}      & $\geq$\,1.24 & array operations, FFT (phonon DOS) \\
\texttt{scipy}      & optional     & Birch--Murnaghan fit in \texttt{plot\_eos.py} \\
\texttt{json}       & stdlib       & snapshot loading \\
\texttt{csv}        & stdlib       & CSV loading \\
\bottomrule
\end{tabular}
\end{center}

No Plotly, Bokeh, VTK, or heavy GUI frameworks.
The suite is intentionally minimal and portable.

% ─────────────────────────────────────────────────────────────────────────────
\section{\S{}G8 --- Implementation Order}

Modules are prioritised by data availability and scientific utility:

\begin{center}
\begin{tabular}{rll}
\toprule
Phase & Module(s) & Blocker \\
\midrule
1 & \texttt{plot\_score\_dist}, \texttt{plot\_score\_radar}         & snapshots already produced \\
1 & \texttt{plot\_lattice\_param}, \texttt{plot\_ecoh\_bar},
    \texttt{plot\_supercell\_scale}                                  & snapshots already produced \\
1 & \texttt{plot\_session\_timeline}                                 & snapshots already produced \\
2 & \texttt{plot\_eos}                                               & EOS CSV engine \\
2 & \texttt{plot\_convergence}                                       & SCF/MD convergence writer \\
3 & \texttt{plot\_rdf}, \texttt{plot\_rdf\_3d}                      & RDF engine \\
3 & \texttt{plot\_thermo}, \texttt{plot\_thermo\_3d}                & MD thermostat writer \\
4 & \texttt{plot\_trajectory\_2d}, \texttt{plot\_trajectory\_3d}    & XYZA MD trajectory writer \\
4 & \texttt{plot\_phonon\_dos}                                       & phonon engine \\
\bottomrule
\end{tabular}
\end{center}

Phase\,1 modules can be written immediately against live snapshot data.
All Phase\,1 modules are 2D; 3D modules appear from Phase\,3 onward.

% ─────────────────────────────────────────────────────────────────────────────
\section{\S{}G9 --- Notes and Constraints}

\begin{itemize}[leftmargin=2em]
  \item No atomistic-level simulation data is invented or hardcoded.
        If a data file does not exist, the module exits with a clear message.
  \item Terminology: use \emph{atomistic} throughout.
        The word \emph{meso} is forbidden in all module code, comments, and axis labels.
  \item All modules must run without a display server when \texttt{--no-anim}
        is passed (\texttt{matplotlib.use('Agg')} fallback).
  \item Figures are research-output quality: axis labels carry units,
        titles carry element symbol and lattice type.
  \item The \texttt{\_style.py} shared module is the single source of truth
        for colours and font settings.
        Individual modules must not hardcode style constants.
\end{itemize}

\vspace{2em}
\hrule
\vspace{0.5em}
{\small\color{dim}%
VSEPR-SIM Graphing Module Plan v1.0 ---
April 2026 ---
Planning document only; no implementation code.
}

\end{document}
